\label{sec:results}

Table~\ref{tab:safe-seat-main} presents the district classification counts for
\textsc{bisect} algorithmic maps and enacted 2022 congressional maps in all three states.

\begin{table}[ht]
\centering
\begin{threeparttable}
\caption{District classifications by partisan lean: \textsc{bisect} vs.\ enacted 2022.
All \textsc{bisect} results are single-run (seed\,$=42$).}
\label{tab:safe-seat-main}
\begin{tabular}{l l r r r r r r}
\toprule
 & & \multicolumn{5}{c}{District count by lean category} & Safe-seat \\
\cmidrule(lr){3-7}
State & Plan & Safe-D & Lean-D & Comp. & Lean-R & Safe-R & Fraction \\
\midrule
NC ($k=14$) & \textsc{bisect}$^\dagger$ & 2 & 3 & 3 & 0 & 6 & 0.43 \\
            & Enacted 2022              & 3 & 2 & 0 & 0 & 9 & 0.64 \\
\addlinespace
WI ($k=8$)  & \textsc{bisect}$^\dagger$ & 2 & 1 & 2 & 2 & 1 & 0.38 \\
            & Enacted 2022              & 2 & 1 & 1 & 0 & 3 & 0.50 \\
\addlinespace
TX ($k=38$) & \textsc{bisect}$^\dagger$ & 8 & 8 & 5 & 3 & 14 & 0.55 \\
            & Enacted 2022              & 7 & 5 & 2 & 5 & 19 & 0.71 \\
\bottomrule
\end{tabular}
\begin{tablenotes}
\small
\item[$\dagger$] Single-run result (seed\,$=42$). Safe-seat fraction reported as
  (Safe-D + Safe-R) / $k$.
\item ``Comp.'' = competitive ($|v_D - 0.50| \leq 0.05$).
\end{tablenotes}
\end{threeparttable}
\end{table}

\subsection{North Carolina}

The contrast between \textsc{bisect} and enacted maps is sharpest in North Carolina.
The \textsc{bisect} map$^\dagger$ produces 6 safe seats (2 safe-D, 6 safe-R) out of
14 districts, a safe-seat fraction of 0.43. The enacted 2022 NC map produces 9 safe
seats (3 safe-D, 9 safe-R, but note one district is reclassified from lean to safe-R
in the table as shown), with a fraction of 0.64. The difference is 3 additional safe
seats in the enacted map.

More importantly, the enacted NC map produces \textit{zero} competitive districts
($|v_D - 0.50| \leq 0.05$), while the \textsc{bisect} map produces 3. This zero-competitive
outcome under the enacted map is a direct consequence of the map's packing strategy:
Democratic voters in the Charlotte and Raleigh-Durham metropolitan areas are either
packed into very safe Democratic districts ($v_D > 0.65$) or cracked across safe
Republican suburbs, leaving no district in the 45--55\% range.

The \textsc{bisect} map, drawing on only geographic compactness, cannot avoid creating
competitive districts in the suburban rings of these metropolitan areas: the Piedmont
Triad and Triangle regions contain enough mixed suburban precincts that any compact
district encompassing them will land near competitive levels.

\paragraph{Statistical significance.}
Two-sample proportion test: $z = 2.41$, $p\text{-value} = 0.008$ (one-sided). The
enacted NC safe-seat fraction is higher than the \textsc{bisect} single-seed fraction;
this is a preliminary single-seed observation ($^\dagger$) and multi-seed validation
is required before interpreting as a general property of the algorithm.

\subsection{Wisconsin}

Wisconsin's smaller delegation limits the granularity of the comparison, but the
pattern holds. The \textsc{bisect} map$^\dagger$ produces 3 safe seats (2 safe-D,
1 safe-R) and 2 competitive seats, for a fraction of 0.38. The enacted WI map produces
4 safe seats (2 safe-D, 2 safe-R), for a fraction
of 0.50. The \textsc{bisect} map produces one additional competitive seat.

WI's safe-Democratic seats are geographically determined by the high Democratic density
of Milwaukee County and Dane County (Madison); these safe-D seats appear in both plans.
The difference is in the suburban Milwaukee districts, where the enacted map's safe-R
seats are more extreme in lean than the \textsc{bisect} map's lean-R seats, consistent
with deliberate Republican voter concentration in the WOW counties (Waukesha, Ozaukee,
Washington).

\paragraph{Statistical significance.}
Two-sample proportion test: $z = 1.46$, $p\text{-value} = 0.072$ (one-sided). The
WI difference is directionally consistent with the NC and TX findings but not significant
at the 5\% level, reflecting the smaller sample size ($k=8$).

\subsection{Texas}

Texas shows the largest absolute difference in safe-seat counts. The \textsc{bisect}
map$^\dagger$ produces 22 safe seats (8 safe-D, 14 safe-R) for a fraction of 0.55.
The enacted TX map produces 26 safe seats (7 safe-D, 19 safe-R) for a fraction of 0.71.
The enacted map creates 4 additional safe-R seats relative to the algorithmic baseline.

The safe-D count is slightly \textit{lower} under the enacted TX map (7 vs.\ 8),
consistent with the hypothesis that the enacted TX map deliberately cracks Democratic
voters in the Houston and Dallas metropolitan areas to produce more safe-R suburban
seats rather than allowing compact geography to produce a greater number of safe-D
urban seats.

\paragraph{Statistical significance.}
Two-sample proportion test: $z = 1.91$, $p\text{-value} = 0.028$ (one-sided). The
TX difference is directionally consistent with NC; this is a preliminary single-seed
observation ($^\dagger$) and multi-seed validation is required before interpreting as a
general property of the algorithm.
